% !TeX root = ../../Book2.tex
% 纲要标题：## 第14章　光滑逼近、Sard 与横截性
% 纲要位置：几何学计划纲要.md，第 1538 行。
% 正文状态：初稿；尚未正式审核。

% 纲要内容开始
% 
% > 本章目标与证明范围：先建立光滑逼近与函数空间的必要拓扑，再完成 Sard 和参数横截性。一般位置的结论具有明确正则性、紧性、相对保持和边界条件。
% 
% 纲要内容结束

\chapter{光滑逼近、Sard 与横截性}
\label{g2:ch:14}
\BookChapterNumber{g2:ch:14}

\begin{chapterabstract}
	方程的解集何时成为流形，以及不规则的相交能否通过微小改变变得规则，是本章的两个基本问题。先以卷积、单位分解和坐标内替换证明相对光滑逼近，说明紧源与非紧源上导数控制的含义；再以秩分层、覆盖估计和维数归纳证明光滑 Sard 定理。横截性将正则值推广到子流形约束，给出逆像的维数、切空间及法丛。参数定理把总体横截化为几乎处处切片横截，局部扰动进一步给出保持闭集附近数据的横截逼近。最后处理边界条件，并由横截同伦构造逆像间的配边。
\end{chapterabstract}

\begin{learninggoals}
	\begin{itemize}
		\item 用局部卷积与截断构造光滑逼近，说明相对条件如何在拼接中保持。
		\item 区分 \(C^0\)、\(C^1\) 及更高阶控制，辨认非紧源上的弱拓扑、强拓扑与紧支改动。
		\item 重建光滑 Sard 定理的证明，并说明临界点集与临界值集的差别及有限可微性的作用。
		\item 将横截条件化为法方向上的满射，计算横截逆像的切空间、维数及法丛。
		\item 从约束流形到参数空间的投影证明参数横截定理，区分稠密、剩余与几乎处处。
		\item 构造相对横截扰动，核对边界限制映射，并用紧性将同伦逆像变成配边。
	\end{itemize}
\end{learninggoals}

考虑平面中的方程 \(y=x^2\) 与直线 \(y=a\)。当 \(a>0\) 时，它们交于两点，每个交点的切线方向都不同；当 \(a<0\) 时没有交点；在 \(a=0\) 时，两点合并成一个相切点。交点数在这个特殊时刻改变，但把参数也加进去，全部解组成的集合 \(\{(a,x):a=x^2\}\) 始终是一条光滑曲线。微分拓扑经常利用这一转换：单个对象可能出现退化，放进一个有足够参数的光滑族后，可以研究哪些切片仍然规则，以及规则切片之间如何联系。

\ProseParagraphBreak

第一卷第11章的正则水平集定理已经说明：当方程组的微分满射时，解集具有预期维数。现在还要回答两个问题。第一，正则值是否足够多，以至于总能作很小的调整选到它？Sard 定理\footnote{Arthur Sard（1909—1980），美国数学家，研究临界值、线性逼近与样条。}给出“临界值为零测集”的答案。第二，约束若是一条曲线或一张曲面，而非一个目标点，什么条件代替满射？所需条件是映射微分的像与约束子流形的切空间共同张成整个目标切空间，即横截性。

\ProseParagraphBreak

本章使用第一卷第11章的微分、逆函数定理和子流形坐标，第12章的局部有限单位分解与光滑截断；具体调用时会恢复条件。分析方面使用紧集上的一致连续性、Taylor 估计以及 Lebesgue 零测集和 Fubini 定理，所需含义在证明处交代。本章的 Sard 与横截性主干针对光滑映射，光滑逼近同时控制任意指定的有限阶导数；最优有限阶 Sard 条件在第14.2节准确陈述并注明原论文的证明位置。有关光滑逼近和参数横截性的参考见\cite[第6章]{Lee2012Smooth}。

% !TeX root = ../../Book2.tex
% 纲要标题：### 14.1 光滑逼近及函数空间
% 纲要位置：几何学计划纲要.md，第 1542 行。
% 正文状态：未撰写；本文件目前是框架。

% 纲要内容开始
% 
% * 局部卷积和单位分解
% * 保持闭子集附近数据的相对逼近
% * 紧源上的 Cʳ 拓扑
% * 非紧源上支集与拓扑的选择
% 
% 纲要内容结束

\section{光滑逼近及函数空间}
\label{g2:sec:1401}

待撰写

\subsection{局部卷积和单位分解}
\label{g2:subsec:1401:01}
待撰写

\subsection{保持闭子集附近数据的相对逼近}
\label{g2:subsec:1401:02}
待撰写

\subsection{\texorpdfstring{紧源上的 \(C^r\) 拓扑}{紧源上的 Cʳ 拓扑}}
\label{g2:subsec:1401:03}
待撰写

\subsection{非紧源上支集与拓扑的选择}
\label{g2:subsec:1401:04}
待撰写
% 纲要条目排为编号小节标题，并进入目录与书签；说明段仍保留于纲要注释中。

% !TeX root = ../../Book2.tex
% 纲要标题：### 14.2 Sard 定理的证明
% 纲要位置：几何学计划纲要.md，第 1549 行。
% 正文状态：初稿；尚未正式审核。

% 纲要内容开始
% 
% * 临界集合的秩分层
% * 秩为零部分的覆盖估计
% * 维数归纳与 Fubini 步骤
% * 正则性阶数及反例范围
% 
% 纲要内容结束

\section{Sard 定理的证明}
\label{g2:sec:1402}

方程 \(f(x)=y\) 的解集能否成为光滑子流形，取决于解上的微分是否满射。单独检查一个指定的 \(y\) 可能很难；Sard 定理给出整体信息：对光滑映射，坏的目标值只占一个零测集。证明必须控制临界点的\emph{像}，因为源中的临界点本身可能很多，甚至充满整个源流形。

\subsection{临界集合的秩分层}
\label{g2:subsec:1402:01}

先明确这里的“坏”。设 \(f:M^m\to N^n\) 是 \(C^1\) 映射，\(N\) 无边界。若 \(\mathrm{d}f_x:\mathrm{T}_xM\to\mathrm{T}_{f(x)}N\) 不满射，称 \(x\) 为\term{临界点}[critical point]；否则称为\term{正则点}[regular point]。临界点的像称为\term{临界值}[critical value]。若 \(f^{-1}(y)\) 中每一点都正则，称 \(y\) 为\term{正则值}[regular value]。特别地，不在像中的点也是正则值；一个值只要有一个临界原像，就已经是临界值。

\ProseParagraphBreak

在 \(n\) 维流形 \(N\) 中，集合 \(E\) 称为\term{零测集}[set of measure zero]，若它在每张坐标图中的像都有 \(n\) 维 Lebesgue 测度零。Euclid 空间中的含义是：给定任意 \(\varepsilon>0\)，能用可数个长方体覆盖它，并使体积总和小于 \(\varepsilon\)。这个条件与坐标无关。在较小的紧坐标块上，光滑坐标变换的导数有界，因而是 Lipschitz 映射；一个半径为 \(r\) 的小球的像落在半径至多为 \(Lr\) 的球中，覆盖总体积至多乘上固定常数。再用可数个这样的坐标块覆盖，便知零测性保持。这里用的是同维坐标变化，无须先在 \(N\) 上选择体积形式。

\begin{example}\label{g2:exm:sard-critical-values}
	映射 \(f:\mathbb R^2\to\mathbb R^2\)，\(f(x,y)=(x,y^2)\)，有
	\[
		\mathrm{D}f_{(x,y)}=\begin{pmatrix}1&0\\0&2y\end{pmatrix}.
	\]
	临界点集是 \(y=0\)，临界值集是目标中的横轴。第二坐标为正的值有两个正则原像，第二坐标为负的值没有原像；这两类值都正则。对于常值映射 \(M\to\mathbb R\)，整个 \(M\) 都由临界点组成，临界值却只有一个。
\end{example}

在坐标域 \(U\subset\mathbb R^m\) 内写 \(C=\{x:\operatorname{rank}\mathrm{D}f_x<n\}\)，并按秩分为
\[
	C^{(q)}=\{x:\operatorname{rank}\mathrm{D}f_x=q\},\qquad 0\le q<n.
\]
这些集合未必都是子流形。分层在这里首先是一种分类：当 \(q\ge1\) 时，某个非零的偏导数让我们把一个输出坐标直接改作输入坐标，进而把问题降低一维。

\begin{lemma}\label{g2:lem:sard-rank-reduction}
	设 \(f:U\to\mathbb R^n\) 光滑，\(U\subset\mathbb R^m\) 开，且 \(\operatorname{rank}\mathrm{D}f_a\ge1\)。重排坐标后，在 \(a\) 的某个邻域上存在光滑源坐标 \((t,z)\in J\times V\subset\mathbb R\times\mathbb R^{m-1}\)，使
	\[
		f(t,z)=(t,G(t,z)),\qquad
		\operatorname{rank}\mathrm{D}f_{(t,z)}=1+\operatorname{rank}\mathrm{D}_zG(t,z).
	\]
\end{lemma}
\begin{proof}
	选择非零偏导数并重排坐标，使 \(\partial f^1/\partial x^1(a)\ne0\)。令
	\[
		\Phi(x)=\bigl(f^1(x),x^2,\ldots,x^m\bigr).
	\]
	其导数行列式正是这个非零偏导数。第一卷第11章第3节的逆函数定理保证 \(\Phi\) 在充分小的邻域上可光滑反解；缩小像邻域，可取乘积盒 \(J\times V\)。在新坐标中，第一分量为 \(t\)，其余分量记作 \(G\)，于是
	\[
		\mathrm{D}f=\begin{pmatrix}1&0\\\partial_tG&\mathrm{D}_zG\end{pmatrix}.
	\]
	利用第一行消去第一列下方的元素，矩阵的秩便等于 \(1+\operatorname{rank}\mathrm{D}_zG\)。
\end{proof}

这里没有要求邻域内的秩恒定；非零偏导数只负责选出一个可用的坐标。真正需要继续处理的，正是其余方向上仍会变化的秩。

\subsection{秩为零部分的覆盖估计}
\label{g2:subsec:1402:02}

若 \(\mathrm{D}f_a=0\)，上面的坐标办法暂时失效。此时 Taylor 展开提供另一种控制：若若干阶导数都为零，一个边长为 \(\delta\) 的源立方体的像就比 \(\delta\) 小得多。源中小立方体的数量约为 \(\delta^{-m}\)，而每个像的 \(n\) 维体积约为 \(\delta^{n(k+1)}\)；当后一个指数更大时，总体积趋于零。

\begin{lemma}\label{g2:lem:sard-flat-cover}
	设 \(U\subset\mathbb R^m\) 开，\(m,n\ge1\)，\(f:U\to\mathbb R^n\) 为 \(C^{k+1}\) 映射，\(k\ge1\)。令
	\[
		A_k=\{x\in U:\partial^\alpha f^i(x)=0\text{，对所有 }1\le|\alpha|\le k,\ 1\le i\le n\}.
	\]
	若 \(n(k+1)>m\)，则 \(f(A_k)\) 为零测集。
\end{lemma}
\begin{proof}
	用可数个包含于 \(U\) 的闭立方体覆盖 \(U\)，只需在其中一个闭立方体 \(Q\) 上证明结论。\(Q\) 上各个 \(k+1\) 阶偏导数有共同的上界。将边长为 \(R\) 的 \(Q\) 均分成 \(s^m\) 个边长 \(R/s\) 的小立方体。对每个与 \(A_k\) 相交的小立方体 \(Q_j\)，选取 \(a_j\in A_k\cap Q_j\)。Taylor 公式的低阶项全部为零，故存在与 \(s,j\) 无关的常数 \(B\)，使
	\[
		|f(x)-f(a_j)|\le B|x-a_j|^{k+1}
		\le B\bigl(\sqrt m R/s\bigr)^{k+1}\qquad(x\in Q_j).
	\]
	因此 \(f(A_k\cap Q)\) 被至多 \(s^m\) 个半径为 \(B(\sqrt m R/s)^{k+1}\) 的球覆盖。若 \(\omega_n\) 表示单位 \(n\) 维球的体积，覆盖总体积至多为
	\[
		s^m\omega_n B^n(\sqrt m R/s)^{n(k+1)}
		=C_Qs^{m-n(k+1)}\longrightarrow0.
	\]
	零测性由此成立；可数并仍为零测集，故得到整个 \(A_k\) 的结论。
\end{proof}

这里的多重指标 \(\alpha=(\alpha_1,\ldots,\alpha_m)\) 记录各坐标求导次数，\(|\alpha|=\alpha_1+\cdots+\alpha_m\)。上式只需普通的有限阶 Taylor 余项估计，不要求函数等于自己的 Taylor 级数。对于光滑函数，在一点的全部导数都为零也不能推出其在邻域内恒定；例如 \(e^{-1/x^2}\) 在 \(x=0\) 补值为零后的函数就具有这种现象。覆盖估计恰好能处理这些极其平坦的点。

\subsection{维数归纳与 Fubini 步骤}
\label{g2:subsec:1402:03}

剩下的工作是把前面两种控制接起来。正秩部分沿一个输出坐标切片；零秩但尚未高度平坦的部分，则落在某个偏导数的正则零集上。两种情形都把源的维数减少一。

\ProseParagraphBreak

切片时使用如下 Fubini 结论：若 \(E\subset\mathbb R\times\mathbb R^{n-1}\) 紧，\(n\ge2\)，且每个截面 \(E_t=\{w:(t,w)\in E\}\) 都为 \((n-1)\) 维零测集，则 \(E\) 为 \(n\) 维零测集。因为紧集是 Borel 可测集，对其示性函数使用非负函数的 Fubini--Tonelli 定理，得到
\[
	\operatorname{vol}_n(E)=\int_{\mathbb R}\operatorname{vol}_{n-1}(E_t)\,\mathrm{d}t=0.
\]
这一紧集切片结论也有只用有限覆盖的证明，见\cite[第126页，引理6.2]{Lee2012Smooth}。下面每次都先限制到紧块，使所处理的像确实可测。

\begin{theorem}[Sard 定理]\label{g2:thm:sard}
	设 \(M^m\) 是有限维、Hausdorff、第二可数的光滑流形，允许有边界；\(N^n\) 是无边界光滑流形，具有相同的分离性与可数性条件。若 \(f:M\to N\) 光滑，则它的临界值集在 \(N\) 中为零测集。因此正则值在 \(N\) 中稠密。
\end{theorem}
\begin{proof}
	先证明无边界情形，并对源维数 \(m\) 归纳。当 \(n=0\) 时，所有微分都满射，临界值集为空。当 \(m=0<n\) 时，第二可数的零维流形是可数离散集，像也可数，故为零测集。以下设 \(m,n\ge1\)，并假设所有源维数小于 \(m\) 的光滑映射均已满足结论。

	取可数的源、目标坐标图，并将源图限制在各目标图的原像内。坐标变化保持微分的秩和目标零测性，故只需讨论 \(f:U\subset\mathbb R^m\to\mathbb R^n\)。除临界集 \(C\) 外，记 \(A_j\) 为前一引理中一至 \(j\) 阶偏导数全为零的集合。它们在 \(U\) 内闭，并满足
	\[
		C\supset A_1\supset A_2\supset\cdots.
	\]
	取整数 \(k\ge1\) 使 \(n(k+1)>m\)。我们依次处理
	\[
		C=(C\setminus A_1)\cup\bigcup_{j=1}^{k-1}(A_j\setminus A_{j+1})\cup A_k.
	\]

	对 \(C\setminus A_1\)，若 \(n=1\)，此集合为空：一维目标的微分不满射就只能为零。若 \(n\ge2\)，\cref{g2:lem:sard-rank-reduction}在每个这样的点附近给出 \(f(t,z)=(t,G(t,z))\)。令 \(G_t(z)=G(t,z)\)。秩公式说明 \((t,z)\) 是 \(f\) 的临界点，当且仅当 \(z\) 是 \(G_t:V\to\mathbb R^{n-1}\) 的临界点。对源坐标盒内任意紧块 \(K\)，集合 \(E=f(C\cap K)\) 紧，而它在目标超平面第一坐标等于 \(t\) 上的截面包含于 \(G_t\) 的临界值集。\(G_t\) 的源维数是 \(m-1\)，由归纳假设，每个这样的截面为零测集；Fubini 结论遂给出 \(E\) 零测。用可数紧块穷尽各坐标盒，再用可数个坐标盒覆盖 \(C\setminus A_1\)，得到 \(f(C\setminus A_1)\) 零测。

	对 \(A_j\setminus A_{j+1}\)，取一点 \(a\)。某个分量的某个 \(j+1\) 阶偏导数在 \(a\) 非零，所以能选一个 \(j\) 阶偏导数 \(h=\partial^\alpha f^i\)，使 \(\mathrm{D}h_a\ne0\)。同时 \(h\) 在 \(A_j\) 上为零。第一卷第11章第3节的隐函数定理说明，在 \(a\) 的充分小邻域 \(V_a\) 内，\(H_a=\{h=0\}\) 是光滑超曲面：选取 \(h\) 的一个非零偏导数，即可把相应坐标解成其余 \(m-1\) 个坐标的光滑函数。于是 \(A_j\cap V_a\subset H_a\)。在这些点上，\(\mathrm{d}f=0\)，故限制映射 \(f|_{H_a}\) 的微分也为零，特别地不满射。归纳假设应用于源维数为 \(m-1\) 的 \(f|_{H_a}\)，说明 \(f(A_j\cap V_a)\) 零测。用可数个这样的邻域覆盖 \(A_j\setminus A_{j+1}\)，便得到这一整部分的像零测。

	最后，\cref{g2:lem:sard-flat-cover}直接给出 \(f(A_k)\) 零测。三部分合起来完成归纳。

	若 \(M\) 有边界，在每张边界坐标图中，光滑映射的坐标表达按定义能在每一点附近延拓到半空间之外的开邻域。原映射在边界点的微分就是这个局部延拓的微分；所以原映射的临界点也是延拓映射的临界点。对可数个局部延拓使用刚证的无边界结论，仍得到临界值集零测。最后，当 \(n\ge1\) 时，任意非空目标开集在某张坐标图中包含一个小球，具有正测度，不能全部由临界值组成；\(n=0\) 时所有值都正则。因此正则值稠密。
\end{proof}

这条证明的分工可由下表记忆。它按照点上的导数选择工具，而没有把整片临界集误当成一个光滑子流形。

\begin{center}
	\begin{tabularx}{\textwidth}{@{}lXX@{}}
		\toprule
		临界点所在部分 & 可用的信息 & 控制临界值的方法 \\
		\midrule
		\(C\setminus A_1\) & 某个一阶偏导数非零 & 一个输出坐标作参数，降维后用 Fubini \\
		\(A_j\setminus A_{j+1}\) & 某个高阶偏导数有正则零集 & 限制到包含该部分的超曲面，降低源维数 \\
		\(A_k\) & 低阶 Taylor 项全部消失 & 小立方体数量与像球体积的覆盖估计 \\
		\bottomrule
	\end{tabularx}
\end{center}

\begin{corollary}\label{g2:cor:sard-lower-dimension}
	设 \(f:M^m\to N^n\) 是上述流形间的光滑映射，且 \(m<n\)。则 \(f(M)\) 是 \(N\) 中的零测集；每个正则值的原像均为空。
\end{corollary}
\begin{proof}
	从 \(m\) 维空间到 \(n\) 维空间的线性映射不可能满射，因此每一点都是临界点。像恰为临界值集，应用\cref{g2:thm:sard}即可。
\end{proof}

例如光滑曲线在平面中的像具有二维测度零，即使曲线在非紧参数区间上的像可能稠密，仍不能因此获得正面积。“稠密”说的是闭包，“零测”说的是体积，二者衡量不同的性质。

\subsection{正则性阶数及反例范围}
\label{g2:subsec:1402:04}

光滑性在上述归纳中使任意阶偏导数都继续光滑。对于有限阶可微映射，\(h=\partial^\alpha f^i\) 的可微阶数会随 \(|\alpha|\) 增大而降低，隐函数定理所得超曲面的正则性也随之降低。因此，那个归纳过程不能直接给出最优的有限阶条件。

\ProseParagraphBreak

有限阶 Sard 定理的准确形式是：对 \(C^r\) 映射 \(f:M^m\to N^n\)，若整数 \(r\ge1\) 满足
\[
	r>\max\{m-n,0\},
\]
则临界值集为零测集。这里 \(M,N\) 可取光滑流形，映射只要求 \(C^r\)，源的边界按局部延拓处理。有限阶版本的证明参见\cite[第885--889页，定理4.1、7.2]{Sard1942}；其中高阶平坦集合的细分估计负责避免上述正则性损失。本章后续参数定理使用已完整证明的光滑版本。

\ProseParagraphBreak

这个阶数条件有可见的几何意义。对平面上的函数，\(C^2\) 足够；单有 \(C^1\) 时，临界值甚至可以覆盖整个区间。下面给出一个具体构造，使“更多可微性”所排除的现象看得见。

\begin{example}\label{g2:exm:sard-c1-counterexample}
	存在 \(C^1\) 函数 \(F:\mathbb R^2\to\mathbb R\)，其临界值集包含 \([0,2]\)。

	从 \([0,1]\) 开始，每一步保留每个区间两端各占其长度 \(\lambda=2/5\) 的闭子区间，得到 Cantor 集 \(E\)。第 \(k\) 步删去 \(2^{k-1}\) 个开间隙，每个长度为 \((1-2\lambda)\lambda^{k-1}\)。固定一个在 \((0,1)\) 上为正、端点处平坦、积分为 \(1\) 的光滑函数 \(\rho\)，并在每个第 \(k\) 步间隙 \(I=(a,b)\) 上放置
	\[
		h(x)=\frac{3^{-k}}{b-a}\rho\left(\frac{x-a}{b-a}\right).
	\]
	在 \(E\) 及 \([0,1]\) 外令 \(h=0\)。间隙上这项的积分为 \(3^{-k}\)，上确界至多为
	\[
		\frac{\|\rho\|_\infty}{3(1-2\lambda)}
		\left(\frac1{3\lambda}\right)^{k-1}\longrightarrow0.
	\]
	这保证 \(h\) 连续：在某个固定间隙的端点附近，连续性来自 \(\rho\) 在端点趋于零；当间隙的级数趋于无穷时，则由上述一致上界趋于零。有限多个低级间隙与余下间隙分别处理，即得到每个 \(E\) 中点处的连续性。

	令 \(g(x)=\int_0^xh(t)\,\mathrm{d}t\)，则 \(g\in C^1(\mathbb R)\)，且 \(g'=0\) 于 \(E\)。全部间隙的积分之和为 \(\sum_{k\ge1}2^{k-1}3^{-k}=1\)。一个第 \(k\) 步保留下来的基本区间，其内各级间隙的积分之和为
	\[
		\sum_{j=k+1}^{\infty}2^{j-k-1}3^{-j}=3^{-k}.
	\]
	所以 \(g\) 把保留区间依次映到标准三分 Cantor 构造中的保留区间，且把每个间隙映到相应的中间三分之一。于是 \(g(E)=C\)，其中 \(C\) 是标准三分 Cantor 集。

	现在取 \(F(x,y)=g(x)+g(y)\)。在 \(E\times E\) 上，\(\mathrm{d}F=(g'(x),g'(y))=0\)，而 \(F(E\times E)=C+C=[0,2]\)。最后这个等式可这样验证：\([-1,1]\) 中每个数都能写成 \(\sum_{j\ge1}2\varepsilon_j3^{-j}\)，其中 \(\varepsilon_j\in\{-1,0,1\}\)，因为第一位所给的三个闭区间覆盖 \([-1,1]\)，选定第一位后再缩放余项即可递归选择。把正、负位分别放到两个只含 \(0,2\) 的三进制展开里，得到 \(C-C=[-1,1]\)。再由 \(1-C=C\)，得到 \(C+C=1+(C-C)=[0,2]\)。
\end{example}

因此，零测结论涉及映射的正则性，而不只是处处存在一个连续的微分。在一般维数中，整数阶阈值的最优性及相应构造见\cite[第890页，第8节]{Sard1942}；对当前 \(m=2,n=1\) 的情形，上例已经完整展示 \(C^1\) 与 \(C^2\) 的差别。另一个边界是源和目标维数：当 \(m<n\) 时，\(C^1\) 已足以让整个像零测：在一个紧坐标立方体上，导数有界使映射 Lipschitz；分割所得的 \(s^m\) 个小立方体各自的像都被半径至多 \(B/s\) 的球包含，故总体积至多为 \(Cs^{m-n}\to0\)。再取可数坐标覆盖即可。

% !TeX root = ../../Book2.tex
% 纲要标题：### 14.3 横截逆像和参数定理
% 纲要位置：几何学计划纲要.md，第 1556 行。
% 正文状态：初稿；尚未正式审核。

% 纲要内容开始
% 
% * 线性横截条件与切空间
% * 逆像的维数及法丛
% * 由总体映射到几乎处处参数
% * 稠密性、剩余性和泛型的区别
% 
% 纲要内容结束

\section{横截逆像和参数定理}
\label{g2:sec:1403}

正则值定理处理方程 \(f(x)=y\)：微分能够独立改变全部输出时，解集有预期维数。现在把一个点 \(y\) 换成子流形 \(S\subset N\)，研究条件 \(f(x)\in S\)。沿着 \(S\) 的输出变化不会破坏这个条件，真正需要控制的是离开 \(S\) 的方向。横截性正是把“全部输出方向”改成“全部法方向”后的满秩条件。

\subsection{线性横截条件与切空间}
\label{g2:subsec:1403:01}

设 \(S\subset N\) 是无边界光滑流形 \(N\) 的嵌入子流形。在 \(y\in S\) 处，用包含映射的微分把 \(\mathrm{T}_yS\) 视为 \(\mathrm{T}_yN\) 的子空间。商空间 \(\mathrm{T}_yN/\mathrm{T}_yS\) 记录法方向：两个向量只要相差一个切于 \(S\) 的向量，就代表同一个法方向。这里还没有选择内积，因而“法”暂时不表示垂直。

\begin{definition}[横截性]\label{g2:def:transversality}
	设 \(M,N\) 为无边界光滑流形，\(S\subset N\) 为无边界嵌入子流形，\(f:M\to N\) 光滑。若对每个 \(x\in f^{-1}(S)\) 都有
	\begin{equation}\label{g2:eq:transverse-linear-condition}
		\mathrm{d}f_x(\mathrm{T}_xM)+\mathrm{T}_{f(x)}S=\mathrm{T}_{f(x)}N,
	\end{equation}
	则称 \(f\) 与 \(S\) \term[hengjiexing]{横截}[transverse]，记作 \(f\pitchfork S\)。若只在 \(x\in K\cap f^{-1}(S)\) 处要求此等式，就说 \(f\) 在 \(K\subset M\) 上与 \(S\) 横截。
\end{definition}

横截条件等价于商映射
\[
	\overline{\mathrm{d}f_x}:\mathrm{T}_xM\longrightarrow\mathrm{T}_{f(x)}N/\mathrm{T}_{f(x)}S,
	\qquad v\longmapsto[\mathrm{d}f_x(v)]
\]
满射。因而只需补足 \(S\) 缺少的方向；\(\mathrm{d}f_x\) 本身可以不满射。当 \(S=\{y\}\) 时，\(\mathrm{T}_yS=0\)，这个条件恰好说 \(y\) 是正则值。淹没与每个嵌入子流形都横截；若 \(f(M)\cap S=\varnothing\)，横截条件也成立，因为没有需要检查的点。记号中的横截并不要求相交。

\begin{example}\label{g2:exm:transverse-curves}
	取 \(S=\mathbb R\times\{0\}\subset\mathbb R^2\)。映射 \(f(t)=(t,t)\) 在唯一交点处给出斜向量 \((1,1)\)，它与 \(S\) 的切向量 \((1,0)\) 张成平面，所以 \(f\pitchfork S\)。映射 \(g(t)=(t,t^2)\) 在零处只给出 \((1,0)\)，与 \(S\) 相切，故不横截。映射 \(h(t)=(t,t^2+1)\) 的像避开 \(S\)，也与 \(S\) 横截。\cref{g2:fig:transversality-curves}将这三种情形并排画出。
\end{example}

\begin{figure}[htbp]
	\centering
	% !TeX root = ../Book2.tex
\begin{tikzpicture}[x=0.9cm,y=0.9cm,>=stealth]
	\begin{scope}
		\draw[->,gray] (-1.6,0)--(1.6,0) node[right] {\(S\)};
		\draw[thick,blue!65!black] (-1.05,-1.05)--(1.05,1.05);
		\fill (0,0) circle (1.4pt);
		\node at (0,-1.6) {横截};
	\end{scope}
	\begin{scope}[xshift=4.5cm]
		\draw[->,gray] (-1.6,0)--(1.6,0) node[right] {\(S\)};
		\draw[thick,blue!65!black,domain=-1.1:1.1,samples=45] plot (\x,{\x*\x});
		\fill (0,0) circle (1.4pt);
		\node at (0,-1.6) {相切};
	\end{scope}
	\begin{scope}[xshift=9cm]
		\draw[->,gray] (-1.6,0)--(1.6,0) node[right] {\(S\)};
		\draw[thick,blue!65!black,domain=-1.0:1.0,samples=45] plot (\x,{\x*\x+0.5});
		\node at (0,-1.6) {空交，也横截};
	\end{scope}
\end{tikzpicture}

	\caption[横截、相切与空交]{横轴表示子流形 \(S\)，曲线表示映射的像。左图在交点处补足法方向；中图的切线仍落在 \(S\) 内；右图没有交点，也满足横截条件。}
	\label{g2:fig:transversality-curves}
\end{figure}

若 \(A,B\subset N\) 都是嵌入子流形，对包含映射 \(i_A:A\hookrightarrow N\) 使用上述定义，就得到两子流形的横截条件
\[
	\mathrm{T}_yA+\mathrm{T}_yB=\mathrm{T}_yN\qquad(y\in A\cap B).
\]
它对 \(A,B\) 对称。两条空间曲线在 \(\mathbb R^3\) 中若有交点，两个一维切空间之和至多二维，所以不可能横截；相交曲线在图上看起来“交叉”，还不足以判断横截性，环境维数也参与条件。

\subsection{逆像的维数及法丛}
\label{g2:subsec:1403:02}

第一卷第11章第4节的嵌入子流形坐标切片说明：若 \(S\) 的余维为 \(c\)，那么在每个 \(y\in S\) 附近，可以找一个淹没 \(q:V\to\mathbb R^c\)，使 \(S\cap V=q^{-1}(0)\)。具体说，选取把 \(S\) 写成最后 \(c\) 个坐标为零的适配坐标，再令 \(q\) 取出这 \(c\) 个坐标。于是 \(f(x)\in S\) 局部变成 \(q(f(x))=0\)。

\begin{theorem}[横截逆像]\label{g2:thm:transverse-preimage}
	设 \(M^m,N^n\) 为无边界光滑流形，\(S\subset N\) 为余维 \(c\) 的无边界嵌入子流形，\(f:M\to N\) 光滑且 \(f\pitchfork S\)。则非空的 \(W=f^{-1}(S)\) 是 \(M\) 的 \((m-c)\) 维嵌入子流形，并且
	\begin{equation}\label{g2:eq:transverse-preimage-tangent}
		\mathrm{T}_xW=\{v\in\mathrm{T}_xM:\mathrm{d}f_x(v)\in\mathrm{T}_{f(x)}S\}.
	\end{equation}
	若 \(c>m\)，则 \(W=\varnothing\)。若 \(S\) 在 \(N\) 中闭，则 \(W\) 在 \(M\) 中闭。
\end{theorem}
\begin{proof}
	在 \(y=f(x)\in S\) 附近取上述 \(q\)。因为 \(\ker\mathrm{d}q_y=\mathrm{T}_yS\)，横截条件等价于 \(\mathrm{d}(q\circ f)_x\) 满射。为看清这一步，任取 \(a\in\mathbb R^c\)，先用 \(q\) 为淹没找到 \(u\in\mathrm{T}_yN\) 满足 \(\mathrm{d}q_y(u)=a\)，再由横截性把 \(u\) 写成 \(\mathrm{d}f_x(v)+w\)，其中 \(w\in\mathrm{T}_yS\)。于是 \(a=\mathrm{d}(q\circ f)_x(v)\)。反过来，若这个复合满射，对任意 \(u\in\mathrm{T}_yN\) 选择 \(v\) 满足 \(\mathrm{d}(q\circ f)_x(v)=\mathrm{d}q_y(u)\)，便有 \(u-\mathrm{d}f_x(v)\in\mathrm{T}_yS\)。

	因此 \(0\) 是 \(q\circ f:f^{-1}(V)\to\mathbb R^c\) 的正则值。第一卷第11章第4节的正则水平集定理给出局部的 \((m-c)\) 维坐标切片，切空间为 \(\ker\mathrm{d}(q\circ f)_x\)，即\eqref{g2:eq:transverse-preimage-tangent}。这些局部切片都取 \(M\) 的子空间拓扑，且坐标转换来自环境坐标的限制，故拼成所述嵌入子流形。若 \(c>m\)，从 \(m\) 维空间到 \(c\) 维空间不可能满射，所以不存在原像点。最后，\(S\) 闭时由 \(f\) 连续得到 \(W\) 闭。
\end{proof}

例如，两张在 \(N^n\) 中横截相交的子流形 \(A^a,B^b\)，其交非空时维数为 \(a+b-n\)，且在交点处
\[
	\mathrm{T}_y(A\cap B)=\mathrm{T}_yA\cap\mathrm{T}_yB.
\]
这是将定理用于 \(i_A:A\hookrightarrow N\) 的直接结果。横截条件保证方程提供了独立约束；一个余维 \(c\) 的目标子流形，恰好在源中施加 \(c\) 个独立方程。

\ProseParagraphBreak

切空间公式还可以记录随点变化的法方向。对嵌入子流形 \(S\subset N\)，定义其\term[facong]{法丛}[normal bundle]为商向量丛
\[
	\nu_N S\coloneqq\mathrm{T}N|_S/\mathrm{T}S.
\]
在适配坐标中，这个商丛由最后 \(c\) 个坐标方向给出局部标架；换坐标诱导光滑可逆的商空间变换，因而确实定义了秩 \(c\) 的光滑向量丛。若选取 Riemann 度量，正交投影将此商丛识别为 \((\mathrm{T}S)^\perp\)，但商丛的定义本身不依赖度量。

\begin{proposition}[逆像的法丛]\label{g2:prop:transverse-normal-bundle}
	在\cref{g2:thm:transverse-preimage}的条件下，令 \(\bar f=f|_W:W\to S\)。微分诱导规范的光滑向量丛同构
	\[
		\nu_M W\longrightarrow\bar f^{\,*}\nu_N S,
		\qquad[v]\longmapsto[\mathrm{d}f(v)].
	\]
	等价地，有短正合列
	\[
		0\longrightarrow\mathrm{T}W\longrightarrow\mathrm{T}M|_W
		\xrightarrow{\ \overline{\mathrm{d}f}\ }\bar f^{\,*}\nu_N S\longrightarrow0.
	\]
\end{proposition}
\begin{proof}
	由\eqref{g2:eq:transverse-preimage-tangent}，\(\overline{\mathrm{d}f_x}\) 的核恰为 \(\mathrm{T}_xW\)，横截性又保证它满射。线性代数的商空间同构定理逐点给出所述同构。它们还须组成光滑丛映射：在源与目标的适配坐标中，微分由光滑矩阵表示，取法方向就是取相应的商矩阵，因而随点光滑。取商后的矩阵为可逆的 \(c\) 阶矩阵，其逆由伴随矩阵除以非零行列式给出，亦光滑。整个定义只使用微分与自然商映射，所以与坐标无关。
\end{proof}

由此可见，横截逆像除了继承预期的维数，也继承目标子流形的法方向。例如，若 \(S\) 的法丛已经给定定向，就可以通过这个同构给 \(W\) 的法丛定向。后面的相交数构造将使用这项信息。

\subsection{由总体映射到几乎处处参数}
\label{g2:subsec:1403:03}

一个指定映射可能不横截。给它增加几个可自由调节的参数后，参数方向有机会补足原来缺少的法方向。例如
\[
	F:\mathbb R\times\mathbb R\longrightarrow\mathbb R,
	\qquad F(a,t)=t^2-a.
\]
固定 \(a=0\) 时，\(F_0(t)=t^2\) 与 \(\{0\}\) 不横截；总体映射却处处是淹没，因为 \(\partial F/\partial a=-1\)。下面的定理把这种可调节性转化为关于单个参数的结论。

\begin{theorem}[参数横截性]\label{g2:thm:parametric-transversality}
	设 \(P,M,N\) 为无边界光滑流形，\(S\subset N\) 为无边界嵌入子流形，\(F:P\times M\to N\) 光滑。若 \(F\pitchfork S\)，则对几乎处处的 \(a\in P\)，切片映射 \(F_a:M\to N\)，\(F_a(x)=F(a,x)\)，满足 \(F_a\pitchfork S\)。这里的例外集在 \(P\) 的每张坐标图内都为 Lebesgue 零测集；特别地，好参数在 \(P\) 中稠密。
\end{theorem}
\begin{proof}
	先把所有满足约束的参数与源点放在一起。由\cref{g2:thm:transverse-preimage}，
	\[
		Z=F^{-1}(S)\subset P\times M
	\]
	是光滑嵌入子流形。考虑投影 \(\pi:Z\to P\)，\(\pi(a,x)=a\)。我们证明：\(a\) 是 \(\pi\) 的正则值，当且仅当 \(F_a\pitchfork S\)。这样就能将横截性问题交给\cref{g2:thm:sard}。

	固定 \((a,x)\in Z\)，令 \(y=F(a,x)\)、\(Q=\mathrm{T}_yN/\mathrm{T}_yS\)。把总体微分送到这个商空间，分别记参数方向与源方向的部分为
	\[
		B:\mathrm{T}_aP\longrightarrow Q,\qquad
		A:\mathrm{T}_xM\longrightarrow Q.
	\]
	于是总体微分为 \((u,v)\mapsto Bu+Av\)。总体横截性表示 \(\operatorname{im}B+\operatorname{im}A=Q\)，而切空间公式给出
	\[
		\mathrm{T}_{(a,x)}Z=\{(u,v):Bu+Av=0\}.
	\]
	投影微分 \(\mathrm{d}\pi_{(a,x)}\) 满射，等价于每个 \(u\) 都能配一个 \(v\)，使 \(Bu+Av=0\)，即 \(\operatorname{im}B\subset\operatorname{im}A\)。结合总体的满射性，这恰好等价于 \(\operatorname{im}A=Q\)，也就是 \(F_a\) 在 \(x\) 处与 \(S\) 横截。对该参数下的所有 \(x\) 同时检查，就得到所需等价；若纤维为空，两边的条件都成立。

	Sard 定理说明 \(\pi\) 的临界值集为零测集，故其补集中的所有参数都横截。任意非空坐标开集都有正测度，不可能整个包含在例外集中，所以好参数稠密。零维参数空间中，Sard 的结论表示没有坏参数，此时上述线性证明也直接给出每个切片横截。
\end{proof}

上述证明使用的不是 \(F\) 的正则值，而是约束流形 \(Z\) 到参数空间的投影的正则值。正则值条件保证：改变参数以后，可以同时改变源点来补偿约束的变化。这正是把总体可调节性传到单个切片的步骤。横截逆像和参数定理也可参看\cite[第6章，第143--146页]{Lee2012Smooth}。

\begin{example}\label{g2:exm:parametric-parabola}
	回到 \(F(a,t)=t^2-a\)，约束流形 \(Z\) 是参数平面中的抛物线 \(a=t^2\)，投影为 \(\pi(t)=t^2\)。它只有一个临界值 \(a=0\)。当 \(a>0\) 时，两个零点 \(\pm\sqrt a\) 处导数非零；当 \(a<0\) 时，没有零点；两类参数都满足横截性。\(a=0\) 是两个零点合并的时刻。参数定理说明坏参数集很小，却并不声称所有参数都好，也不保证不同好参数的原像具有相同的连通分支数。
\end{example}

\subsection{稠密性、剩余性和泛型的区别}
\label{g2:subsec:1403:04}

参数定理使用有限维参数空间的测度；“映射经过任意小扰动就能横截”则使用函数空间的拓扑。要把这两种表述衔接起来，先需要一个稳定性事实：已经在紧集上满足的横截条件，可以承受足够小的一阶扰动。

\begin{proposition}[紧集上的横截开放性]\label{g2:prop:transversality-open-on-compact}
	设 \(M,N\) 为无边界光滑流形，\(S\subset N\) 为闭的无边界嵌入子流形，\(K\subset M\) 紧。若光滑映射 \(f:M\to N\) 在 \(K\) 上与 \(S\) 横截，则 \(f\) 有一个弱 \(C^1\) 邻域，其中每个光滑映射在 \(K\) 上仍与 \(S\) 横截。特别地，当 \(M\) 紧时，全部 \(C^1\) 横截映射构成 \(C^1(M,N)\) 中的开集。
\end{proposition}
\begin{proof}
	对每个 \(x\in K\)，区分两种情形。若 \(f(x)\notin S\)，因 \(S\) 闭，可取 \(x\) 的相对紧邻域 \(U_x\)，使 \(f(\overline U_x)\) 落在 \(N\setminus S\) 中；充分小的 \(C^0\) 扰动在较小闭邻域上仍避开 \(S\)。这一控制可用覆盖该紧像的有限张坐标图表达。

	若 \(f(x)\in S\)，取适配坐标及局部定义函数 \(q:V\to\mathbb R^c\)。横截性使 \(\mathrm{D}(q\circ f)\) 的某个 \(c\) 阶子式在 \(x\) 非零。缩小相对紧的 \(U_x\)，可使 \(f(\overline U_x)\subset V\)，且这个子式的绝对值在 \(\overline U_x\) 上有正下界。对 \(g\) 在此紧集上的值与一阶导数施加足够小的坐标误差，就能保持 \(g(\overline U_x)\subset V\) 并使同一个子式仍非零。因此 \(q\circ g\) 在这里为淹没，尤其在 \(g^{-1}(S)\cap U_x\) 上满足横截条件。

	这些邻域覆盖 \(K\)。取有限子覆盖，再同时施加相应的有限个紧集 \(C^1\) 条件，便得到所需弱 \(C^1\) 邻域。取 \(K=M\) 得最后结论；该证明只涉及一阶导数，因而也适用于 \(C^1\) 映射。
\end{proof}

闭性在这里用来阻止新的交点从 \(S\) 的遗漏极限点附近出现。例如 \(S=\{1,1/2,1/3,\ldots\}\subset\mathbb R\) 是非闭的零维嵌入子流形。从单点流形到 \(\mathbb R\) 的映射 \(f(*)=0\) 避开 \(S\)，因而横截；任意接近它的某个常值映射 \(f_j(*)=1/j\) 却不横截。这个例子即使源紧，开放性仍然失效。

\ProseParagraphBreak

源的紧性与所选拓扑也须匹配。映射 \(f(x)=x\) 从 \(\mathbb R\) 到自身与 \(\{0\}\) 横截。取光滑函数 \(\rho\)，使其在零附近等于 \(1\)，在 \([-1,1]\) 外为零，令
\[
	g_j(x)=x\bigl(1-\rho(x-j)\bigr).
\]
在每个固定紧集上，\(g_j\) 最终与 \(f\) 完全相同，所以 \(g_j\to f\) 于弱 \(C^\infty\) 拓扑；但 \(g_j\) 在 \(j\) 附近恒为零，均不横截。紧集上的开放性无法一次控制逃往无穷远处的交点。

\begin{definition}[稠密与剩余]\label{g2:def:generic-residual}
	在拓扑空间 \(X\) 中，若 \(A\) 与每个非空开集相交，则称 \(A\) \term[choumi]{稠密}[dense]。若 \(A\) 包含一个可数个开稠密集的交，则称 \(A\) 为\term[shengyuji]{剩余集}[residual set]。一个性质在剩余集上成立时，通常称其为\term[fanxing]{泛型}[generic]性质，使用时须同时说明空间及其拓扑。
\end{definition}

稠密性保证任意小的邻域里有好对象，剩余性则记录可数个开稠密要求能够同时成立。在 Baire 空间中，每个可数开稠密交都稠密。例如 \(\mathbb R\) 具有这个性质：给定非空开区间及开稠密集 \(U_1,U_2,\ldots\)，依次在前一个区间内部与 \(U_j\) 的交中选非空闭区间，令长度趋零；实数的完备性给出一个属于所有区间的点。这个嵌套区间论证说明了完备性如何使可数条件同时实现。一般拓扑空间未必有此性质：在 \(\mathbb Q\) 的相对拓扑中，枚举其点为 \(q_1,q_2,\ldots\)，每个 \(\mathbb Q\setminus\{q_j\}\) 都开且稠密，交却为空。

\ProseParagraphBreak

对光滑映射空间，本章使用的具体结论是：若 \(M,N\) 无边界，\(S\subset N\) 是闭的无边界嵌入子流形，则
\[
	\mathcal T_S=\{f\in C^\infty(M,N):f\pitchfork S\}
\]
在弱 \(C^\infty\) 拓扑及\secref{g2:sec:1401}定义的强 \(C^\infty\) 拓扑中都稠密且为剩余集。稠密性由下一节的\cref{g2:thm:relative-transversality}证明：其中取固定闭集为空，便能在任意给定的强邻域内得到横截映射。剩余性的推导则很短。取紧集穷尽 \(K_1\subset K_2\subset\cdots\)，\(\bigcup_jK_j=M\)，令
\[
	\mathcal U_j=\{f:f\text{ 在 }K_j\text{ 上与 }S\text{ 横截}\}.
\]
由\cref{g2:prop:transversality-open-on-compact}，每个 \(\mathcal U_j\) 在这两种拓扑中都开；全局横截映射的稠密性又使每个 \(\mathcal U_j\) 稠密。最后 \(\mathcal T_S=\bigcap_j\mathcal U_j\)。这里稠密性有单独的扰动证明，不需要先对整个流形映射空间证明一个 Baire 定理。

\ProseParagraphBreak

“几乎处处”与“泛型”也不是同一个量词。将有理数枚举为 \(q_i\)，对每个 \(k\) 取以这些点为中心的开区间，使总长度小于 \(2^{-k}\)，其并 \(U_k\) 开且稠密。于是 \(G=\bigcap_kU_k\) 是稠密剩余集，却有测度零，因为它包含在每个总长度小于 \(2^{-k}\) 的 \(U_k\) 中。相应地，\(\mathbb R\setminus G\) 为满测集，却是可数个闭且无内点的集合之并。函数空间中的泛型结论不附带一个未指定的概率分布；参数横截定理的几乎处处结论则明确使用有限维参数坐标中的 Lebesgue 测度。

% !TeX root = ../../Book2.tex
% 纲要标题：### 14.4 相对与带边界横截性
% 纲要位置：几何学计划纲要.md，第 1563 行。
% 正文状态：初稿；尚未正式审核。

% 纲要内容开始
% 
% * 局部扰动和有限覆盖
% * 保持指定区域不变
% * 边界限制映射的横截条件
% * 同伦族与配边的生成
% 
% ---
% 
% 纲要内容结束

\section{相对与带边界横截性}
\label{g2:sec:1404}

参数定理要求先有一个总体横截的映射族。对任意给定映射，怎样造出这样的族？在目标坐标中加一个小向量，就能补足缺少的方向；在源端乘截断函数，又能把改变限制在指定区域。将这个局部操作逐块实施，便得到保持已给数据的一般位置结果。边界问题还需检查沿边界本身是否有足够的方向，最后才能把一族逆像拼成配边。

\subsection{局部扰动和有限覆盖}
\label{g2:subsec:1404:01}

\begin{lemma}[坐标内的横截扰动]\label{g2:lem:local-transverse-perturbation}
	设 \(M,N\) 是无边界光滑流形，\(S\subset N\) 是嵌入子流形，\(f:M\to N\) 光滑。设 \(K\subset U\subset M\)，其中 \(K\) 紧，\(U\) 开，\(f(U)\) 落在一张目标坐标图内。则可在 \(U\) 内选择紧支扰动，使所得光滑映射 \(g\) 在 \(K\) 上横截于 \(S\)，并在 \(M\setminus U\) 上等于 \(f\)。扰动可以在任意预先指定的有限阶 \(C^r\) 意义下任意小。
\end{lemma}
\begin{proof}
	记目标坐标为 \(y:V\to\mathbb R^n\)。第一卷第12章第3节的光滑截断结论保证存在 \(\rho:M\to[0,1]\)，在 \(K\) 的某邻域上为 \(1\)，且支集紧含于 \(U\)。取充分小的参数球 \(B\subset\mathbb R^n\)，定义
	\[
		F(a,x)=\begin{cases}
			y^{-1}\bigl(y(f(x))+\rho(x)a\bigr),&x\in U,\\
			f(x),&x\notin\operatorname{supp}\rho.
		\end{cases}
	\]
	两个表达在交叠处一致；\(f(\operatorname{supp}\rho)\) 紧含于 \(V\)，故可统一缩小 \(B\)，使平移后的点仍在坐标像内。在 \(\rho>0\) 的区域，关于参数的微分在坐标中为 \(\rho(x)I_n\)，所以总体映射 \(F\) 是淹没。\cref{g2:thm:parametric-transversality}应用于 \(B\times\{\rho>0\}\)，说明几乎每个任意小的 \(a\) 都使 \(F_a\) 在这里横截于 \(S\)，尤其在 \(K\) 上横截。

	令 \(g=F_a\)。在固定紧支集上，\(F_a\) 的各阶坐标导数连续依赖于 \(a\)，并在 \(a=0\) 时等于 \(f\) 的相应导数。因此让 \(a\) 足够小，即可满足任意给定的有限阶误差界。对非紧源上的强拓扑，紧支集只与局部有限控制族中的有限多个紧块相交，同一结论仍适用。
\end{proof}

这个构造说明参数为什么应在目标坐标中取值：源端原本可能缺少的法向方向，由 \(a\) 直接补上。截断函数的过渡区不必横截；下一步会用覆盖把所有尚未处理的点纳入某个 \(\rho=1\) 的区域。

\ProseParagraphBreak

先设 \(M\) 紧，\(S\) 闭。取有限个紧块 \(K_i\subset U_i\)，其内部覆盖 \(M\)，每个 \(f(U_i)\) 落在一张目标图中。依次应用引理，使第 \(i\) 步在 \(K_i\) 上横截。已经处理的 \(K_1\cup\cdots\cup K_{i-1}\) 是紧集；\cref{g2:prop:transversality-open-on-compact}保证，只要本步足够小，就保留那里的横截性。再把总误差分配成有限份，最后得到任意接近 \(f\) 的处处横截映射。这里保持的是先前取得的满秩条件，先前的函数值可以继续改变。

\subsection{保持指定区域不变}
\label{g2:subsec:1404:02}

若映射在某个区域已经用于拼接，那里常常必须逐点不变。例如同伦的两个端点已经给定，不能在“取得一般位置”时顺便更换端点。相对横截性把这种要求直接列入结论。

\begin{theorem}[相对横截逼近]\label{g2:thm:relative-transversality}
	设 \(M,N\) 是有限维、Hausdorff、第二可数的无边界光滑流形，\(S\subset N\) 是闭嵌入子流形，\(A\subset M\) 闭。若光滑映射 \(f:M\to N\) 在 \(A\) 的一个开邻域上横截于 \(S\)，则对每个有限整数 \(r\ge1\) 及 \(f\) 的每个强 \(C^r\) 邻域 \(\mathcal U\)，存在光滑映射 \(g\in\mathcal U\)，使 \(g\pitchfork S\)，并使 \(g=f\) 在 \(A\) 的某个开邻域上成立。\(g\) 还可选为与 \(f\) 光滑同伦，且同伦在这个较小邻域上不变。\(A=\varnothing\) 时即得到普通横截逼近。
\end{theorem}
\begin{proof}
	设 \(f\) 在开集 \(U_0\supset A\) 上横截。由第一卷第4章第3—4节的正规性与覆盖收缩，选开集 \(W,V\)，使
	\[
		A\subset W\subset\overline W\subset V\subset\overline V\subset U_0.
	\]
	可先缩小 \(\mathcal U\)，使其中每个映射在 \(\overline V\) 上横截。具体做法是在 \(\overline V\) 附近取局部有限的紧坐标块，在每块上用紧集横截开放性给出正的 \(C^1\) 容差；强拓扑允许同时要求这些逐块不等式。因而接下来只需处理 \(M\setminus V\)。

	对这个闭集取可数、在整个 \(M\) 上局部有限的相对紧开集族 \(U_i\)，每个 \(\overline{U_i}\) 避开 \(\overline W\)，且 \(f(\overline{U_i})\) 含于一个目标坐标图。取紧块 \(K_i\subset U_i\)，使其内部覆盖 \(M\setminus V\)。这里使用的是第一卷第12章第3节构造单位分解时的局部有限收缩；尤其局部有限性在 \(W\) 的边缘也成立。用逐块 \(C^0\) 容差保证后续映射仍留在所选目标图内。

	从 \(f_0=f\) 开始，在第 \(i\) 步应用\cref{g2:lem:local-transverse-perturbation}，得到 \(f_i\)，使其在 \(K_i\) 上横截，改动支集紧含于 \(U_i\)。同时让改动足够小，以保留 \(K_1,\ldots,K_{i-1}\) 上的横截性，并满足原邻域的误差要求。下面说明无穷多步中如何保留这些要求，避免只在每个有限阶段正确。

	在固定的源、目标坐标块中用各阶导数的上确界测量差。若原邻域在第 \(\ell\) 块给出容差 \(\varepsilon_\ell\)，令第 \(i\) 步在它碰到的各块上耗用不到 \(2^{-i-1}\varepsilon_\ell\)。第 \(i\) 步完成后，紧集开放性给出 \(\delta_i>0\)：以后只要在覆盖 \(K_i\) 的有限坐标块上与 \(f_i\) 的 \(C^1\) 差小于 \(\delta_i\)，便仍横截。对每个 \(j>i\)，额外要求第 \(j\) 次改动在这些块上的大小小于 \(2^{-(j-i+1)}\delta_i\)。总和小于 \(\delta_i\)。一次改动有紧支，所遇原控制块有限，先前的 \(K_i\) 也只有有限多个；所以每一步只是对参数添加有限多个严格上界，可以同时满足。

	每一点的某邻域只遇到有限多个 \(U_i\)，因此 \(f_i\) 在该邻域最终不再改变，定义出光滑映射 \(g\)。上述误差和保证 \(g\in\mathcal U\) 并在每个 \(K_i\) 上横截；在 \(\overline V\) 上的横截性由最初缩小的邻域保证。这些集合覆盖 \(M\)，故 \(g\pitchfork S\)。所有改动支集避开 \(\overline W\)，于是 \(g=f\) 于 \(W\)。

	最后，在各次坐标平移中把选定参数 \(a_i\) 换成 \(t a_i\)，按相同次序组合，\(0\le t\le1\)。各步可一并选得对所有 \(t\) 满足前述坐标内的误差预算，所以组合有定义。源端的局部有限性保证每点附近只有有限次组合，从而给出联合光滑的 \(H(t,x)\)，满足 \(H(0,x)=f(x)\)、\(H(1,x)=g(x)\)，且在 \(W\) 上恒等于 \(f\)。这个同伦本身的每张切片不必横截。
\end{proof}

\begin{corollary}[余维较大的集合可以避开]\label{g2:cor:avoid-high-codimension}
	设 \(M,N\) 是无边界光滑流形，\(S\subset N\) 是闭嵌入子流形，且 \(\dim M<\operatorname{codim}_N S\)。每个光滑映射 \(M\to N\) 都可任意光滑逼近为像避开 \(S\) 的映射，并保持同伦类；若原映射在闭集 \(A\) 的某邻域上避开 \(S\)，可在较小邻域内保持它不变。
\end{corollary}
\begin{proof}
	相对横截逼近给出 \(g\pitchfork S\)。若 \(g(x)\in S\)，则 \(\mathrm{T}_xM\to\mathrm{T}_{g(x)}N/\mathrm{T}_{g(x)}S\) 必须满射，与维数不等式矛盾。因此 \(g^{-1}(S)=\varnothing\)。任意 \(C^\infty\) 邻域只需满足构成它的一组有限阶基本条件，选足够大的有限 \(r\) 即可。
\end{proof}

例如三维空间中的一条闭路可以作任意小的扰动避开给定点；对一张曲面也成立，因为点的余维为3。若障碍物是一张曲面，余维只有1，上述维数论证便不再给出避开结论。相交数将在后续章节解释何时交点反而不能全部消除。

\subsection{边界限制映射的横截条件}
\label{g2:subsec:1404:03}

\begin{theorem}[带边界的横截逆像]\label{g2:thm:boundary-transverse-preimage}
	设 \(M^m\) 是光滑带边界流形，\(N\) 是无边界光滑流形，\(S\subset N\) 是无边界、余维为 \(c\) 的嵌入子流形。若光滑映射 \(f:M\to N\) 以及限制 \(f|_{\partial M}\) 都横截于 \(S\)，则 \(Z=f^{-1}(S)\) 若非空，就是 \(m-c\) 维带边界嵌入子流形，且
	\[
		\partial Z=Z\cap\partial M,\qquad
		\mathrm{T}_xZ=(\mathrm{d}f_x)^{-1}(\mathrm{T}_{f(x)}S).
	\]
	在 \(x\in\partial Z\) 处还有
	\[
		\mathrm{T}_x(\partial Z)=\mathrm{T}_xZ\cap\mathrm{T}_x(\partial M),\qquad
		\mathrm{T}_xZ+\mathrm{T}_x(\partial M)=\mathrm{T}_xM.
	\]
\end{theorem}
\begin{proof}
	内部点由\cref{g2:thm:transverse-preimage}处理。边界点处取 \(S\) 的局部定义函数 \(q:V\to\mathbb R^c\)，即 \(S\cap V=q^{-1}(0)\)，且 \(\ker\mathrm{d}q=\mathrm{T}S\) 沿 \(S\) 成立。令 \(h=q\circ f\)。边界横截条件等价于 \(\mathrm{d}(h|_{\partial M})\) 满射，这正是第一卷第11章第4节“带边界的正则水平集”所需的条件。

	为看清边界为何保留下来，在半空间坐标 \((x^1,\ldots,x^m)\)、\(x^m\ge0\) 中，满射性使前 \(m-1\) 个方向中的某 \(c\) 列给出可逆子式。重排这些坐标后，以
	\[
		\bigl(h^1(x),\ldots,h^c(x),x^{c+1},\ldots,x^{m-1},x^m\bigr)
	\]
	作新坐标。逆函数定理适用于坐标函数的局部光滑延拓，且最后一个坐标仍为 \(x^m\)，故局部逆保持半空间。\(Z\) 在这些坐标中就是前 \(c\) 项为零的半空间切片。切片边界、切空间及最后一个向内方向分别给出陈述中的各个等式。包含映射在这些图中为坐标嵌入，图册与子空间拓扑相容。
\end{proof}

\begin{example}\label{g2:exm:boundary-transversality-failure}
	在上半平面 \(M=\{(x,y):y\ge0\}\) 上令 \(f(x,y)=y+x^2\)，取 \(S=\{0\}\subset\mathbb R\)。由于 \(\partial_y f=1\)，\(f\) 处处为淹没；然而 \(f^{-1}(0)=\{(0,0)\}\)，并不是预期的一维带边界流形。缺少的条件是 \(f|_{\partial M}(x)=x^2\) 在零点处的横截性。改成 \(g(x,y)=x\)，逆像为向上的半直线，才具有一维内部及一个边界点。完整切空间含向内法向，而边界限制只能使用沿边界的方向，两项检查因此有不同作用。
\end{example}

相对逼近也有相同的边界版本：在\cref{g2:thm:relative-transversality}中允许源 \(M\) 带边界，要求 \(f\) 与 \(f|_{\partial M}\) 在固定区域都横截，则可使最终 \(g\) 与 \(g|_{\partial M}\) 同时横截。证明中的坐标扰动在半空间上仍有意义；在 \(\rho>0\) 处，总体映射限制到 \(B\times\operatorname{Int}M\) 和 \(B\times\partial M\) 都因参数方向而成为淹没。对这两张无边界源的映射分别应用参数定理，避开两个零测异常集的并，即可同时满足条件。边界上的满射还蕴含完整切空间上的满射。紧块开放性分别应用于内部的完整导数和边界切向导数，此后的有限覆盖或局部有限构造完全沿用上述误差选择。

\subsection{同伦族与配边的生成}
\label{g2:subsec:1404:04}

一条同伦使每个点的像随时间移动；若整条同伦横截于 \(S\)，所有相交点的运动轨迹就组成高一维的流形。它的边界记录起点和终点的逆像。这是由一般位置构造几何不变量的基本方式。

\begin{definition}[配边]\label{g2:def:unoriented-cobordism}
	设 \(Z_0,Z_1\) 是紧无边界 \(k\) 维光滑流形。若存在紧 \((k+1)\) 维光滑带边界流形 \(W\)，以及微分同胚 \(\partial W\cong Z_0\amalg Z_1\)，就称 \(Z_0,Z_1\)\term[peibian]{配边}[cobordant]。此处暂不指定定向；有向版本还需规定边界的诱导定向。
\end{definition}

\begin{theorem}[同伦的横截逆像]\label{g2:thm:transverse-homotopy-cobordism}
	设 \(M^m\) 是紧无边界光滑流形，\(N\) 是无边界光滑流形，\(S\subset N\) 是闭、余维 \(c\le m\) 的嵌入子流形。若两个光滑映射 \(f_0,f_1:M\to N\) 均横截于 \(S\)，且连续同伦，则存在光滑同伦 \(H:[0,1]\times M\to N\)，在两个端点的时间邻域分别恒等于 \(f_0,f_1\)，并且 \(H\pitchfork S\)。于是
	\[
		W=H^{-1}(S),\qquad
		\partial W=\bigl(\{0\}\times f_0^{-1}(S)\bigr)
		\amalg\bigl(\{1\}\times f_1^{-1}(S)\bigr)
	\]
	给出两个逆像之间的配边。
\end{theorem}
\begin{proof}
	从一条连续同伦 \(F\) 出发，取光滑非降函数 \(\lambda:[0,1]\to[0,1]\)，在0附近等于0、在1附近等于1。把 \(F\) 换成 \(F(\lambda(t),x)\)，所得同伦在两端时间邻域中等于光滑的 \(f_0,f_1\)。\cref{g2:thm:relative-smooth-approximation}适用于这个带边界源，给出保持较小端点邻域的光滑同伦 \(F'\)。

	在这些端点邻域内，\(F'\) 以及它的边界限制都横截于 \(S\)，因为 \(f_0,f_1\) 已横截。应用刚证明的带边界相对横截逼近，固定再小一些的端点邻域，得到所需的 \(H\)。\cref{g2:thm:boundary-transverse-preimage}给出 \(W\) 的流形结构及边界公式。最后，\(S\) 闭，所以 \(W\) 是紧空间 \([0,1]\times M\) 的闭子集，因而紧；这一点正是配边定义中的紧性要求。
\end{proof}

\begin{example}\label{g2:exm:zeros-appear-in-pairs}
	在 \(S^1\) 上考虑 \(f_a(e^{\mathrm i\theta})=\cos\theta-a\)。当 \(|a|<1\) 时，零点有两个且都是正则零点；当 \(|a|>1\) 时没有零点。沿 \(a\) 从0变到2的总体函数，其参数导数等于 \(-1\)，故总体逆像是光滑一维流形，虽然 \(a=1\) 的单张切片有临界零点。两个零点的轨迹在该临界时刻接合并消失。这说明总体横截足以构造配边，却不要求每个时刻的切片都横截。
\end{example}

一般地，当 \(m=c\) 时，横截逆像是有限点集。上述配边是紧一维流形，其边界点数为偶数。可以直接用有限图解释这件事：取有限条闭坐标弧，其相对内部覆盖 \(W\)，遇到边界点时使用半区间图。把所有弧端点和 \(\partial W\) 放入有限集 \(P\)。任一坐标弧删去 \(P\) 后的各开段，在 \(W\setminus P\) 中既开又闭，故其连通分支就是有限条这样的开段。这把 \(W\) 写成有限图：内部顶点的度数为2，边界顶点的度数为1。由各边恰有两个端点，全部顶点度数之和为边数的两倍，故度数为1的顶点有偶数个。因此 \(f_0^{-1}(S)\) 与 \(f_1^{-1}(S)\) 的点数模2相同。后续相交数理论将在加入定向后，把这个奇偶计数细化为带符号计数。


\begin{chaptersummary}
	局部卷积提供平滑操作，截断与局部有限覆盖将它拼成保持闭集附近数据的整体逼近。导数的控制阶数决定哪些性质稳定；非紧源上还要区分逐紧集控制与处处给定容差。光滑 Sard 定理把临界集分为可降维的部分和高阶平坦的部分，分别以归纳、Fubini 和覆盖体积估计处理。横截性要求微分在约束的法方向上满射，因此逆像具有预期维数，法丛则是目标法丛的拉回。参数横截定理将切片的横截性等价为约束流形投影的正则值条件；局部坐标扰动使这一结论可以用于任意映射，并保持指定区域。源带边界时，边界限制映射须另外横截，逆像的边界才恰是端点或边界数据的逆像。对于紧源与闭约束，横截同伦的逆像给出配边，进而保持零维逆像点数的奇偶性。
\end{chaptersummary}

\BookExercises

\begin{exercise}\label{g2:ex:cr-convergence-orders}
	对整数 \(k\ge0\)，在 \(S^1\) 上定义 \(h_n(e^{\mathrm{i}\theta})=n^{-k-1}\sin(n\theta)\)。判断 \(h_n\) 在哪些有限阶 \(C^r\) 拓扑中趋于零，并说明为何它不在 \(C^\infty\) 拓扑中趋于零。
\end{exercise}
\begin{solution}
	对第 \(j\) 阶角变量导数，其最大绝对值为 \(n^{j-k-1}\)。故当 \(r\le k\) 时，所有 \(j\le r\) 的导数一致趋于零，得到 \(C^r\) 收敛。第 \(k+1\) 阶导数的最大绝对值恒为 \(1\)，所以 \(r\ge k+1\) 时不收敛。圆周可用有限个角坐标块覆盖，且换角坐标只差常数，以上计算确实检验了圆周的 \(C^r\) 拓扑。\(C^\infty\) 收敛要求每个有限阶都收敛，第 \(k+1\) 阶已经排除了它。
\end{solution}

\begin{exercise}\label{g2:ex:sphere-valued-approximation}
	设 \(M\) 为紧光滑流形，\(f:M\to S^q\subset\mathbb R^{q+1}\) 连续，\(q\ge1\)。若光滑 \(u:M\to\mathbb R^{q+1}\) 满足 \(\sup_M|u-f|<1/2\)，证明 \(g=u/|u|\) 为光滑映射 \(M\to S^q\)，且 \(g\) 与 \(f\) 同伦。给出 \(\sup_M|g-f|\) 的误差上界，并说明若 \(u=f\) 于一个区域，同伦在该区域有何性质。
\end{exercise}
\begin{solution}
	因为 \(|f|=1\)，有 \(|u|\ge1-|u-f|>1/2\)，所以归一化有定义且光滑。又
	\[
		|g-f|\le|g-u|+|u-f|=\bigl|1-|u|\bigr|+|u-f|\le2|u-f|,
	\]
	故 \(\sup_M|g-f|\le2\sup_M|u-f|\)。令 \(v_t=(1-t)f+tu\)，则 \(|v_t-f|=t|u-f|<1/2\)，于是 \(v_t\ne0\)，而 \(H(x,t)=v_t(x)/|v_t(x)|\) 是从 \(f\) 到 \(g\) 的连续同伦。若 \(u=f\) 于某区域，那里 \(v_t=f\) 且 \(|f|=1\)，故整个同伦固定不动。
\end{solution}

\begin{exercise}\label{g2:ex:uniform-not-strong-convergence}
	取非负 \(b\in C^\infty_c(\mathbb R)\)，使 \(b(0)=1\)、\(\operatorname{supp}b\subset[-1/4,1/4]\)，并令 \(u_n(x)=n^{-1}b(x-n)\)。证明 \(u_n\) 在每个有限阶的全局一致导数范数下趋于零，但不在强 \(C^0\) 拓扑中趋于零。
\end{exercise}
\begin{solution}
	对固定 \(r\)，有
	\[
		\max_{0\le j\le r}\sup_{x\in\mathbb R}|u_n^{(j)}(x)|
		=\frac1n\max_{0\le j\le r}\sup_{x\in\mathbb R}|b^{(j)}(x)|\longrightarrow0.
	\]
	强拓扑允许位置相关的误差。在局部有限紧块 \(K_m=[m-1/2,m+1/2]\)、\(m\in\mathbb Z\) 上同时要求 \(\sup_{K_m}|v|<e^{-|m|}/2\)，得到零映射的一个强 \(C^0\) 邻域。然而 \(u_n(n)=1/n>e^{-n}/2\)，故没有任何 \(u_n\) 属于这个邻域。这也表明非紧源的强拓扑比全局一致拓扑更细。
\end{solution}

\begin{exercise}\label{g2:ex:critical-collapse}
	设 \(f:\mathbb R^2\to\mathbb R^2\)，\(f(x,y)=(x,xy)\)。求临界点集、临界值集及所有正则值的原像。说明临界点集的维数与临界值集的维数为何不同。
\end{exercise}
\begin{solution}
	导数矩阵为 \(\begin{pmatrix}1&0\\y&x\end{pmatrix}\)，行列式为 \(x\)，故临界点集为直线 \(\{0\}\times\mathbb R\)。这条直线全部被映到 \((0,0)\)，所以临界值集只有 \(\{(0,0)\}\)。对于 \(a\ne0\)，值 \((a,b)\) 的原像为唯一的 \((a,b/a)\)，它正则；对于 \(a=0,b\ne0\)，原像为空，也正则。唯一的临界值 \((0,0)\) 的原像是整条临界直线。这里一个一维临界点集被压成零维的临界值集，Sard 定理控制的是后者在目标中的测度。
\end{solution}

\begin{exercise}\label{g2:ex:regular-values-not-open}
	令 \(f:\mathbb R\to\mathbb R\)，\(f(x)=e^{-x}\sin x\)。证明 \(0\) 是正则值，但正则值集不是开集。再证明：若 \(M,N\) 为无边界光滑流形、\(f:M\to N\) 光滑且固有，则正则值集是开集。这里固有指紧集的原像紧。
\end{exercise}
\begin{solution}
	\(f'(x)=e^{-x}(\cos x-\sin x)\)，故临界点恰为 \(x_k=\pi/4+k\pi\)，\(k\in\mathbb Z\)，对应临界值为 \((-1)^ke^{-x_k}/\sqrt2\)。这些值均不为零，并在 \(k\to+\infty\) 时趋于零。另一方面，\(f^{-1}(0)=\{k\pi:k\in\mathbb Z\}\)，各点导数为 \((-1)^ke^{-k\pi}\ne0\)，故 \(0\) 正则，而它的任意邻域都含临界值。

	一般情形下，临界点集 \(C\subset M\) 闭，因为满秩矩阵构成开集。设临界值序列 \(y_j=f(x_j)\) 收敛到 \(y\)，其中 \(x_j\in C\)。取 \(y\) 的一个相对紧坐标邻域；去掉有限项后，\(y_j\) 和 \(y\) 位于某个紧坐标块 \(K\) 内。固有性给出 \(f^{-1}(K)\) 紧，故 \(x_j\) 有子列收敛到 \(x\in C\)，且连续性给出 \(f(x)=y\)。于是临界值集闭，其补集即正则值集开。首个例子不固有：任意包含 \(0\) 的小闭区间的原像都包含充分大的正实数，因而不紧。
\end{solution}

\begin{exercise}\label{g2:ex:c1-small-dimensional-image}
	设 \(U\subset\mathbb R^m\) 开，\(f:U\to\mathbb R^n\) 为 \(C^1\) 映射，且 \(m<n\)。不使用有限阶 Sard 定理，直接证明 \(f(U)\) 为 \(n\) 维零测集，并推广到 \(C^1\) 流形映射。
\end{exercise}
\begin{solution}
	先在一个闭立方体 \(Q\subset U\) 上工作。导数有界给出 \(Q\) 上的 Lipschitz 常数 \(L\)。将边长为 \(R\) 的 \(Q\) 分成 \(s^m\) 个小立方体，每个像均包含在半径为 \(L\sqrt mR/s\) 的球中，故 \(f(Q)\) 的覆盖总体积至多为
	\[
		s^m\omega_n(L\sqrt mR/s)^n=C_Qs^{m-n}\longrightarrow0.
	\]
	用可数个这样的闭立方体覆盖 \(U\)，即得结论。对于流形，取可数坐标图并将源图限制到目标图的原像，使坐标表达成为上述 \(C^1\) 映射。坐标图之间的 \(C^1\) 变换在紧块上 Lipschitz，保持同维零测集，再取可数并即可。源有边界时，在半立方体上使用局部 \(C^1\) 延拓，得到同样估计。
\end{solution}

\begin{exercise}\label{g2:ex:transversality-saddle-graph}
	设 \(f:\mathbb R^2\to\mathbb R^3\)，\(f(x,y)=(x,y,x^2-y^2)\)。分别对竖直轴 \(L=\{(0,0,z):z\in\mathbb R\}\) 及水平面 \(P_a=\{z=a\}\) 判断 \(f\) 的横截性，并写出原像的维数及切空间。
\end{exercise}
\begin{solution}
	与 \(L\) 相交只发生在 \((x,y)=(0,0)\)。此处 \(\mathrm{d}f\) 的像为水平的 \(xy\) 平面，与 \(\mathrm{T}_0L\) 的竖直方向共同张成 \(\mathbb R^3\)，故 \(f\pitchfork L\)。原像为单点，维数 \(2-2=0\)，切空间为零。

	\(P_a\) 的法方向由第三坐标给出，复合函数为 \(h(x,y)=x^2-y^2\)，微分为 \(\mathrm{d}h_{(x,y)}(u,v)=2xu-2yv\)。其唯一临界点为原点，临界值为零。因此 \(a\ne0\) 时 \(f\pitchfork P_a\)，原像为一维双曲线，切空间为 \(\{(u,v):xu-yv=0\}\)。\(a=0\) 时原像为两条直线 \(y=x\) 与 \(y=-x\) 的并，原点处不横截，该交叉在原点附近不是一维流形；离开原点后仍有上述一维切空间。
\end{solution}

\begin{exercise}\label{g2:ex:parametric-sufficient-not-necessary}
	设 \(F:\mathbb R^2\to\mathbb R^2\)，\(F(a,t)=(t,t^3-at)\)，目标子流形为横轴 \(S\)。求出哪些 \(a\) 使 \(F_a\pitchfork S\)，并检查总体映射 \(F\) 是否与 \(S\) 横截。这说明参数定理的总体假设是何种条件？
\end{exercise}
\begin{solution}
	切片的交点满足 \(t(t^2-a)=0\)，法方向导数为 \(3t^2-a\)。当 \(a<0\) 时只有 \(t=0\) 一个交点，导数为 \(-a\ne0\)；当 \(a>0\) 时有 \(0,\pm\sqrt a\) 三个交点，对应导数为 \(-a,2a,2a\)，均不为零。当 \(a=0\) 时只有零点，导数也为零。因此好参数恰为 \(\mathbb R\setminus\{0\}\)。

	总体映射的法分量为 \(H(a,t)=t^3-at\)，其微分为 \(-t\,\mathrm{d}a+(3t^2-a)\,\mathrm{d}t\)，在 \((0,0)\) 处为零，且此点映入 \(S\)。故 \(F\) 不横截。几乎处处参数横截仍然成立，说明总体横截性是参数定理给出的充分条件，并非其结论成立的必要条件。
\end{solution}

\begin{exercise}\label{g2:ex:mobius-regular-zero}
	把 Möbius 线丛的总空间写成 \(E=(\mathbb R\times\mathbb R)/((t,v)\sim(t+1,-v))\)，零截面记为 \(S\)。证明不存在光滑函数 \(h:E\to\mathbb R\)，同时满足 \(h^{-1}(0)=S\) 与 \(0\) 为正则值。说明此现象与法丛有什么关系。
\end{exercise}
\begin{solution}
	若这样的 \(h\) 存在，令 \(H(t,v)=h([t,v])\)。则 \(H(t,0)=0\)，故 \(\partial H/\partial t(t,0)=0\)。正则值条件要求 \(b(t)=\partial H/\partial v(t,0)\) 处处非零。商关系给出 \(H(t+1,-v)=H(t,v)\)，对 \(v\) 求导得 \(b(t+1)=-b(t)\)。连续函数 \(b\) 在任何 \([t,t+1]\) 的两端异号，由介值定理必有零点，矛盾。

	从法丛看，若 \(h\pitchfork\{0\}\)，\cref{g2:prop:transverse-normal-bundle}会给出 \(\nu_E S\cong S\times\mathbb R\)。但零截面的法方向是纤维方向，沿底圆周绕一周发生符号翻转，正是 Möbius 线丛本身。上述 \(b(t)\) 的符号论证也直接证明了这个法线丛不能平凡化。
\end{solution}

\begin{exercise}\label{g2:ex:finite-simultaneous-transversality}
	设 \(M\) 为紧无边界光滑流形，\(N\) 为无边界光滑流形，\(S_1,\ldots,S_k\subset N\) 为闭的无边界嵌入子流形。在 \(C^\infty(M,N)\) 的 \(C^\infty\) 拓扑中，证明同时与所有 \(S_i\) 横截的映射构成开稠密集。能否直接由各次横截性推出映射也与所有非空交集 \(S_i\cap S_j\) 横截？
\end{exercise}
\begin{solution}
	由\cref{g2:prop:transversality-open-on-compact}及\cref{g2:thm:relative-transversality}，每个集合 \(\mathcal T_i=\{f:f\pitchfork S_i\}\) 都开且稠密。其有限交开。给定任意非空开集 \(U\)，先与 \(\mathcal T_1\) 相交，所得仍为非空开集；再依次与其余 \(\mathcal T_i\) 相交，每次由稠密性保留非空性。故有限交也稠密。

	最后的推断不成立，即使交集本身是子流形也如此。取 \(M=S^1\)、\(N=\mathbb R^2\)，\(f(\theta)=(\sin\theta,\sin\theta)\)，令 \(S_1,S_2\) 为两条坐标轴。两类交点都在 \(\theta=0,\pi\)，导数为 \((\cos\theta,\cos\theta)\)，与每条轴分别张成平面，所以两次横截性均成立。但是 \(S_1\cap S_2=\{0\}\)，从一维切空间到二维切空间的微分不可能满射，故 \(f\) 不与该交集横截。同时满足各个子流形的横截条件，并不等于对其交集施加了独立的全部法向约束。
\end{solution}

\begin{exercise}\label{g2:ex:boundary-parabola-levels}
	在上半平面 \(M=\{(x,y):y\ge0\}\) 上令 \(f_a(x,y)=y+x^2-a\)。分别判断 \(f_a\) 与 \(f_a|_{\partial M}\) 何时横截于 \(\{0\}\)，求 \(f_a^{-1}(0)\) 及其边界，并解释 \(a=0\) 时何处失去预期维数。
\end{exercise}
\begin{solution}
	完整微分为 \(2x\,\mathrm{d}x+\mathrm{d}y\)，处处满射，故所有 \(a\) 都使 \(f_a\pitchfork\{0\}\)。边界限制为 \(x^2-a\)：若 \(a>0\)，零点为 \(\pm\sqrt a\)，导数分别为 \(\pm2\sqrt a\)，因而横截；若 \(a<0\)，没有零点，也横截；若 \(a=0\)，零点处导数为零，不横截。

	当 \(a>0\) 时，逆像是弧 \(\{(x,a-x^2):-\sqrt a\le x\le\sqrt a\}\)，以 \(x\) 为参数得到一维带边界流形，边界是 \((\pm\sqrt a,0)\)。当 \(a<0\) 时逆像为空。当 \(a=0\) 时逆像只有原点，维数为零，而无边界正则值公式会给出 \(2-1=1\)。此时约束将唯一的向内方向消耗掉，边界切向方向却无法补偿，正是附加边界条件失效的位置。
\end{solution}

\begin{exercise}\label{g2:ex:relative-value-versus-neighborhood}
	令 \(f(x)=x^2\)、\(A=\{0\}\subset\mathbb R\)、\(S=\{0\}\subset\mathbb R\)。证明不存在与 \(f\) 在 \(A\) 的某邻域相同且横截于 \(S\) 的光滑映射。若只要求 \(g(0)=f(0)\)，构造在每个固定紧区间上任意 \(C^r\) 接近 \(f\) 且横截于 \(S\) 的光滑映射，\(r<\infty\)。
\end{exercise}
\begin{solution}
	若 \(g=f\) 于零点的某邻域，则 \(g(0)=0\)、\(g'(0)=0\)，所以 \(g\) 在这个原像点不横截。只固定点值时，取 \(g_\varepsilon(x)=x^2+\varepsilon x\)，\(\varepsilon\ne0\)。它有两个零点 \(0,-\varepsilon\)，其导数分别为 \(\varepsilon,-\varepsilon\)，均非零，故处处横截于 \(S\)。在 \([-R,R]\) 上，差的零阶上界为 \(|\varepsilon|R\)，一阶上界为 \(|\varepsilon|\)，更高阶差为零，故所有有限阶都随 \(\varepsilon\to0\) 逼近。保持邻域内的映射会同时固定导数，只固定闭集上的值则未必如此。
\end{solution}

\begin{exercise}\label{g2:ex:circle-regular-zero-parity}
	设 \(h:S^1\to\mathbb R\) 光滑，且 \(0\) 为正则值。用横截同伦的逆像证明 \(h^{-1}(0)\) 有偶数个点。再对每个整数 \(k\ge0\) 给出恰有 \(2k\) 个正则零点的例子。
\end{exercise}
\begin{solution}
	正则零点孤立，而零点集是紧圆周的闭子集，故它为有限点集。\(h\) 与常值函数 \(1\) 通过 \((1-t)h+t\) 连续同伦，两端映射都横截于 \(\{0\}\)。\cref{g2:thm:transverse-homotopy-cobordism}给出一个紧一维流形，其边界为 \(h^{-1}(0)\) 与空集的不交并。正文的有限图度数论证说明其边界点数为偶数，故结论成立。

	当 \(k=0\) 时取常值函数 \(1\)。当 \(k\ge1\) 时取 \(h(e^{\mathrm i\theta})=\sin(k\theta)\)。零点为 \(\theta=j\pi/k\)、\(0\le j<2k\)，各点导数为 \(k(-1)^j\ne0\)，恰得 \(2k\) 个正则零点。
\end{solution}

\begin{exercise}\label{g2:ex:zeros-escape-to-infinity}
	取 \(M=\mathbb R\)，考虑 \(H:[0,1]\times M\to\mathbb R\)、\(H(t,x)=e^x-t\)。证明总体映射及其边界限制都横截于 \(\{0\}\)，求 \(H^{-1}(0)\)。解释为何两端的正则零点数奇偶性不同，以及这为何不违背本章的配边结论。
\end{exercise}
\begin{solution}
	\(\mathrm{d}H=-\mathrm{d}t+e^x\,\mathrm{d}x\) 处处满射。\(t=0\) 时函数 \(e^x\) 没有零点，横截条件自动成立；\(t=1\) 时只有 \(x=0\) 一个零点，导数为 \(1\)，也横截。总体逆像为
	\[
		W=\{(t,\log t):0<t\le1\},
	\]
	它是与 \((0,1]\) 微分同胚的一维带边界流形，边界只有 \((1,0)\)。当 \(t\downarrow0\) 时，\(x=\log t\to-\infty\)，所以 \(W\) 非紧，缺失的端点逃向源的无穷远处。两端分别有0个和1个零点；配边结论要求源紧，正是为了保证逆像紧并排除这种现象。
\end{solution}
